\section[Meshing]{\Gls{meshing}}
Let now \glssymbol{Omega} denote the whole \gls{domain} and divide it into \gls{N} elements
\begin{subequations} \label{eqn:Element_Locations}
	\begin{align}
		\glssymbol{Omega} &= \left[ 0, 1 \right] = \sum \limits_{i = 1}^{\gls{N}} \gls{Omegai}, \label{eqn:Whole_Region}	\\
		\gls{Omegai} &= \left[ \gls{xi},  x_{i + 1}, \right], \quad i = 1, 2, \dots, \gls{N}.	\label{eqn:Element_No_i}
	\end{align}
\end{subequations}
%
This introduces the following concepts:
\begin{definition}
	\Glspl{nodePoint} are distinct and fixed points in the \gls{domain}, i.e.,
	\begin{align}
		x_{i - 1} < & x_{i} < x_{i + 1}, \quad i = 2, 3, \dots, \gls{N}, \\
		x_1  & = 0, \\
		x_{N + 1} & = 1.
	\end{align}
\end{definition}
\begin{definition}
	\Glspl{element} are distinct parts of the \gls{domain}, i.e.,
	\begin{equation}
		\left\{ \glsentryname{Omegae}: \forall x \vert \quad x_{e} \le x \le x_{e + 1}, \quad e = 1, 2, \dots, \gls{N} \right\} .
	\end{equation}
\end{definition}
Note that element boundaries and locations are expressed in \gls{globalCoord}.
\begin{definition}
	\Gls{meshing} is the process of dividing the region into smaller pieces, called elements in FEM terminology. This process
	determines the number of elements, \gls{N}, and the locations of the $N+1$ nodes. Another term for meshing is discretization.
\end{definition}
\begin{remark}
	Two adjacent elements, $\Omega_e$ and $\Omega_{e+1}$, have exactly one common point, $x_{e+1}$.
\end{remark}
\begin{definition}
	A mesh is uniform if all elements have the same size. In the present one-dimensional case, the nodes are therefore equidistant.
\end{definition}
\begin{remark}
	A common procedure is to create an initial mesh, solve the problem, analyze the solution, and then refine or coarsen the mesh
	in regions where the solution varies rapidly or slowly, respectively. This iterative process, \gls{adaptivemeshing}, produces
	a \gls{nonUniformMeshing} and can improve solution accuracy.
\end{remark}
\begin{figure}
	\includegraphics[width = \textwidth]{../Figures/Meshing}
	\caption{A Typical \Gls{region} with \Glspl{element} and \Glspl{nodePoint}}
	\label{fig:Meshing}
\end{figure}
Consider now an approximation to the \gls{exactSol} given by
\begin{equation} \label{eqn:Truncated_FEM_Solution}
	\glssymbol{uM} = \sum_{i = 1}^{\gls{M}} \glsentryname{ai} \gls{phii}.
\end{equation}
This is a series spanned by a total of \gls{M} linearly independent functions \gls{phii}, also called \glspl{shapeFunction} or
\glspl{interpolationFunction}. The series is in addition truncated because we have replaced infinity ($\infty$) on the upper bound of
summation by \gls{M}. The solution in \cref{eqn:Truncated_FEM_Solution} is thus a linear combination of $M$ \glspl{basisFunction}.
and the task is then to determine the \gls{M} coefficients $a_j$, such that \cref{eqn:FEM_Problem}, as well as
\glspl{boundaryCondition} in \cref{eqn:FEM_Left_Boundary,eqn:FEM_Right_Boundary} are satisfied. In FEM-terminology we say that
the model has $M$ \glspl{dof}. \Glspl{basisFunction} have typically \gls{localSupport}, i.e., they are defined over a few \glspl{element}
and are identically $0$ otherwise. \\

The terms \Glsfirst{finite} and \Glsfirstplural{element} in \gls{FEM} refer to the facts that the solution is a truncated series and the
\gls{region} is divided into subregions as explained above. From now on we denote the \gls{FEMSol} just by $u  = u \left( x \right)$
and drop the truncation reference. The first derivative of $u$ may be written as
\begin{equation} \label{eqn:u_Prime}
	\frac{du}{dx} = \sum_{i = 1}^{\gls{M}} a_i \glsentryname{phiiPrime}.
\end{equation}
A particularly interesting choice of weight functions is \glspl{basisFunction}, i.e.,
\begin{subequations} \label{eqn:Weight_Function}
	\begin{align}
		\gls{wi} &= \gls{phii}, & i &= 1, 2, \dots, \gls{M},
		\shortintertext{such that}
		\glsentryname{wiPrime} &= \glsentryname{phiiPrime} & i &= 1, 2, \dots, \gls{M},
	\end{align}
\end{subequations}
Let us for simplicity assume a \gls{uniformMeshing}. Length of all elements is then equal and all nodes are
equidistant Dividing the whole \gls{region} $x \in \left[ 0, 1 \right] $ into \gls{N} \glspl{element} will then result in the
following element locations and element length
\begin{subequations}
	\begin{align}
		\glssymbol{Omega} &= \left[ 0, 1 \right] = \sum \limits_{i = 0}^{\gls{N}} \gls{Omegai}, \\
		\gls{h} & = \frac{x_{N + 1} - x_{1}}{\gls{N}} = \frac{1 - 0}{\gls{N}} = \frac{1}{\gls{N}}.
	\end{align}
	where
	\begin{alignat}{2}
		\gls{Omegai}	&= \left[ \gls{xi},  x_{i + 1} \right], \quad && i = 1, 2, \dots, \gls{N}, \\
		\gls{xi} 		&= \left( i - 1 \right) h, 				\quad && i = 1, 2, \dots, N + 1.
	\end{alignat}
\end{subequations}
 %
